\documentclass[11pt,a4paper]{article}
\usepackage[margin=1in]{geometry}
\usepackage[utf8]{inputenc}
\usepackage[english]{babel}
\usepackage{amsmath,amssymb}
\usepackage{microtype}
\usepackage{hyperref}

\setlength{\parindent}{0pt}
\setlength{\parskip}{0.75\baselineskip}
\sloppy

\title{Elementary Particle as a Closed Wave of Field 01\\\large Scalar Profile, Normal-Retention Interpretation, Mass, Memory, and the Absence of a Point Singularity}
\author{Nikolai Dereviankin}
\date{}

\newcommand{\field}{\mathcal{F}_{01}}
\newcommand{\phase}{\varphi}
\newcommand{\normal}{\mathcal{N}}
\newcommand{\memory}{\mathcal{M}}
\newcommand{\contour}{\mathcal{C}}

\begin{document}
\maketitle

\begin{center}
\textit{Working version v0.3 --- English research draft for public discussion}
\end{center}

\begin{abstract}
This paper formulates a first core element of Field 01: an elementary particle is interpreted not as a material point, but as a stable closed wave configuration of a field. Two regimes of the field---fixation $0$ and direction $1$---form a rhythm from which phase circulation, a local scalar / VEV-like profile, its Field 01 normal-retention interpretation, mass, and memory can be described. A photon is interpreted as an open transport of phase $1$, whereas a massive elementary particle appears only when the circulation closes and a local scalar / VEV-like profile is retained. This approach replaces the point-like core by a dynamical core and thereby avoids an internal point singularity in the model description of an elementary particle. A black-hole horizon appears as a limiting regime of the same logic: suppression of the local scalar / VEV-like profile and transition of phase memory into a surface record.
\end{abstract}

\textbf{Keywords:} Field 01; closed wave; elementary particle; phase circulation; scalar profile; VEV-like profile; normal-retention interpretation; mass; memory; quantum field theory; horizon; holography.

\section{How to Read This Paper}

This paper is a working research text, not a completed physical theory. Its purpose is to make a first step from the book-level intuitive description of Field 01 toward a formal language that can be discussed, criticized, and gradually compared with established physics.

The reader should distinguish three levels of statements:
\begin{enumerate}
\item \textbf{Established physics}: quantum fields, vacuum correlations, photon energy $E=\hbar\omega$, mass--energy equivalence, the stress--energy tensor, and horizon entropy.
\item \textbf{Interpretation proposed by Field 01}: a particle is described as a stable field regime, a photon as open phase transport, and a massive particle as a closed node with a local scalar / VEV-like profile that receives a normal-retention interpretation in Field 01.
\item \textbf{Open hypothesis}: the scalar / VEV-like profile, its normal-retention interpretation, memory, phase tension, horizon surface recording, and the non-literal evaporation picture used in Field 01 require further mathematical formalization.
\end{enumerate}

Therefore, all formulas not derived from a strict action principle or from a standard theory should be read as working notation and as directions for future formalization. Field 01 does not claim that the Standard Model is wrong, nor does it replace quantum field theory. It offers an additional phase-topological language whose consistency must be tested.

\section{Status of This Paper in the Research Program}

At the present checkpoint, the particle paper is the most developed paper in the Field 01 public archive. Other directions---cosmology, horizon interpretation, gravitational language, information, and possible extensions toward Standard Model language---remain in progress and should not be read as completed conclusions.

This order is methodological. If the minimal object cannot be defined carefully---a particle as a stable field regime---then larger constructions lose their basis. The central test question of the present version is therefore not the whole Field 01 program, but the narrower question:
\begin{center}
\textit{can an elementary particle be described as a stable closed phase configuration of a field?}
\end{center}

The public repository and open discussions are intended to test this question. Supportive interpretations are welcome, but criticism is equally important: mathematical contradictions, comparison with known soliton and topological models, incompatibility with quantum field theory, missing references, and suggestions for an action, energy functional, or stability condition.

This paper should therefore be read as the central working node around which the other texts are being organized. It does not reduce the value of the Russian book-level presentation; instead, it translates its intuitive core into a form that can be checked mathematically and discussed in English.

\section{Aim of the Paper}

Field 01 proposes to treat physical structures not as a collection of ready-made objects in a pre-given space, but as stable regimes of a single field. This paper isolates the minimal core of the model: a possible description of an elementary particle as a closed wave.

The main aim is to introduce a working vocabulary that can connect the intuitive description with a mathematical framework. The proposed formalization is not presented as a replacement for the Standard Model. It is an interpretive and hypothetical layer that may help organize questions about particles, mass, memory, and horizons.

The central thesis is:
\begin{quote}
An elementary particle can be interpreted as a stable closed phase configuration of Field 01, in which the phase circulation is self-consistent and a local scalar / VEV-like profile is retained; the language of normal retention is an interpretation of that scalar-profile layer, not an already established new standard degree of freedom.
\end{quote}

\section{Two Regimes of Field 01}

Let $\field$ denote the field considered in the model. At the qualitative level it is described by two elementary regimes:
\begin{itemize}
\item $0$ --- fixation, holding, boundary, recording;
\item $1$ --- direction, motion, circulation, phase transfer.
\end{itemize}

These symbols should not be read as binary digits in an ordinary computational sense. They are minimal names for two complementary modes of field behavior. The regime $1$ expresses directional phase motion, while the regime $0$ expresses the ability of the field to hold, close, or record a configuration.

A physical process in Field 01 is not the motion of an object through an empty background. It is a change of regime in the field itself. The basic rhythm can be written schematically as
\begin{equation}
0 \leftrightarrow 1.
\end{equation}

This expression is not yet an equation of motion. It is a compact notation for the alternation between fixation and directed phase change.

\section{Phase and Circulation}

To describe a wave-like regime we introduce a phase variable
\begin{equation}
\phase = \phase(x,t),
\end{equation}
where $x$ denotes position in an effective space and $t$ denotes an effective time parameter. At this stage both space and time are treated as descriptive variables, not as independently explained fundamental entities of the model.

A phase circulation along a contour $\contour$ may be written as
\begin{equation}
\oint_{\contour} d\phase.
\end{equation}

If the circulation does not close, the phase is transported without forming a localized node. If the circulation becomes self-consistent, a closed wave configuration can arise. A minimal topological condition for such a configuration is
\begin{equation}
\oint_{\contour} d\phase = 2\pi n, \qquad n \in \mathbb{Z}.
\end{equation}

This condition is not meant to define a complete particle spectrum. It only expresses the idea that a stable node requires a quantized or self-consistent phase return.

\section{Minimal Mathematical Setup}

To move from image to testable structure, it is useful to write a minimal set of variables. Let the local field state be approximated by an amplitude and a phase,
\begin{equation}
\psi(x)=A(x)e^{i\phase(x)}.
\end{equation}
For a closed node it is also useful to consider two phase branches,
\begin{equation}
\psi_a(x)=A_a(x)e^{i\phase_a(x)}, \qquad a=1,2,
\end{equation}
and their local separation,
\begin{equation}
\Delta\phase(x)=\phase_1(x)-\phase_2(x).
\end{equation}
These two branches are not two substances. They are a way to describe internal self-consistency of circulation: the same wave structure must be able to return to itself after a complete phase cycle.

In this notation a closed particle is not a point. It is a map from a circulation contour into a state space,
\begin{equation}
\Gamma:S^1\longrightarrow\mathcal{S}_{01},
\end{equation}
where $S^1$ denotes the closed circulation parameter and $\mathcal{S}_{01}$ is the still-to-be-defined state space of Field 01. The winding number can then be written as
\begin{equation}
n=\frac{1}{2\pi}\oint_{\contour}d\phase.
\end{equation}

For physical interpretation, winding alone is not enough. At least four conditions are required:
\begin{enumerate}
\item \textbf{closure}: observable quantities return to themselves after a full circulation;
\item \textbf{localization}: the node energy is concentrated in a finite region and approaches a vacuum-like background outside it;
\item \textbf{regularity}: the energy density does not contain a point infinity;
\item \textbf{stability}: a small deformation does not immediately destroy the node or lower its energy without a barrier.
\end{enumerate}

These conditions are the first mathematical filter for the model. If they cannot be realized in an explicit energy functional or action, then ``particle as a closed wave'' remains only a metaphor. If they can be realized, the next question is whether the resulting stable classes can be compared with known particle masses, spins, and charges.

\section{Photon as Open Circulation}

In standard physics a photon is a quantum of the electromagnetic field, with energy
\begin{equation}
E = \hbar\omega.
\end{equation}

In the language of Field 01 this relation can be read as a link between energy and the frequency of phase rhythm. A photon is interpreted as an open transport of the regime $1$: phase is carried, but no local scalar / VEV-like profile is retained.

Schematically,
\begin{equation}
\text{photon}: \qquad \normal = 0, \qquad \text{open phase transport}.
\end{equation}

This is not intended to replace quantum electrodynamics. It is a model-specific interpretation of why the photon has no rest mass: in Field 01, mass is associated with the retention of a closed phase structure and a local scalar / VEV-like profile, while the photon remains an open propagating mode.

\section{Elementary Particle as a Closed Wave}

An elementary particle is interpreted as a stable closed configuration of the field. Unlike a photon, it is not merely a transfer of phase along an open direction. It is a self-returning phase regime.

The schematic condition for such a regime is
\begin{equation}
\oint_{\contour} d\phase = 2\pi n.
\end{equation}

Here $\contour$ does not have to be understood as a material wire or a literal circular path. It denotes the effective phase cycle that returns the field to a self-consistent state.

In this interpretation the particle is not a small ball and not a mathematical point with an infinite core. It is a localized mode of field rhythm. Its stability is produced by closure, not by a hard inner substance.

\section{Scalar Profile and the Interpretation of Local Depth}

A closed phase circulation introduces an additional feature absent in the open photon mode: a local scalar / VEV-like profile. The primary object here is a scalar profile $N$ or VEV-like profile. The symbol $\normal$ is kept as an interpretive Field 01 label for local depth or normal retention:
\begin{equation}
N=N[\Delta\phase], \qquad \normal\equiv\text{interpretation}(N).
\end{equation}

This scalar / VEV-like profile is not a material rod inside the particle. In Field 01 language it is a way of describing local depth associated with the coupling of phase branches. The normal-retention language is interpretation-layer vocabulary, not an already established independent standard field.

If the phase separation disappears,
\begin{equation}
\Delta\phase\to 0,
\end{equation}
then the scalar profile is suppressed,
\begin{equation}
N\to 0, \qquad \normal\to 0\;\text{as an interpretive limit}.
\end{equation}

Thus the distinction between a photon and a massive particle can be expressed as
\begin{equation}
\begin{aligned}
\text{photon}:& \qquad N = 0, \\
\text{massive particle}:& \qquad N \neq 0.
\end{aligned}
\end{equation}

This is a working hypothesis of the model. A strict version would require a geometric definition of $N$ and a dynamical equation for its stability.

\section{Localization and Absence of a Point Singularity}

An important requirement for a particle as a closed wave is finite energy. If $\varepsilon(x)$ denotes an effective energy density, a localized node should satisfy
\begin{equation}
E_{\mathrm{node}}=\int_{\mathbb{R}^3}\left(\varepsilon(x)-\varepsilon_{\mathrm{vac}}\right)d^3x<\infty.
\end{equation}
Here $\varepsilon_{\mathrm{vac}}$ denotes the background relative to which the configuration energy is measured. This does not define the final Lagrangian, but it fixes an obligatory requirement: particle mass should not arise from infinite self-energy concentrated at a point center.

In a spherically symmetric approximation, regularity of the core can be expressed as
\begin{equation}
\varepsilon(0)<\infty, \qquad \int_0^{r_0}4\pi r^2\varepsilon(r)\,dr<\infty.
\end{equation}
If the local scalar / VEV-like profile is written as $N(r)$, natural regular boundary conditions are
\begin{equation}
N(0)=N_c<\infty, \qquad \partial_rN(0)=0,
\end{equation}
while outside the node one may require
\begin{equation}
N(r)\to N_{\mathrm{vac}} \quad \text{or} \quad N(r)\to 0,
\end{equation}
depending on the future formalization.

In this form, absence of a point singularity becomes a mathematical condition rather than only a philosophical thesis. The model must produce finite energy, a smooth or weakly singular field configuration, and no need to place mass at a point of zero size. If a chosen action does not satisfy these conditions, that version of the model must be rejected or changed.

\section{Mass as Energy of Retaining the Scalar Profile}

In standard physics mass is related to energy by
\begin{equation}
E = mc^2,
\end{equation}
and, in the Standard Model, elementary particle masses are connected with their coupling to the Higgs field. Field 01 does not deny this structure. It asks whether mass can also be described in an additional phase-geometric language.

In the model, mass is interpreted as the energy cost of maintaining a closed phase configuration and its local scalar / VEV-like profile. A schematic expression is
\begin{equation}
m c^2 \sim E_{\mathrm{closed}}[\phase,N].
\end{equation}

A possible functional form is
\begin{equation}
E_{\mathrm{closed}} = \int_{\Sigma} \left( A |\nabla \phase|^2 + B |N|^2 + V(\phase,N) \right) d\Sigma,
\end{equation}
where $A$, $B$, and $V$ are not yet fixed from a fundamental action. The purpose of this expression is only to show the direction of formalization: mass should correspond to the energy of a stable closed solution, not to an independent substance placed inside the particle.

In words:
\begin{quote}
Mass is associated with the energy required to hold a closed phase rhythm and prevent it from unfolding into an open mode.
\end{quote}

\section{Stability of the Closed Node}

The next level of testing is stability. A closed phase circulation must not only be drawn; it must arise as a stationary configuration of some energy functional. In minimal form this can be written as
\begin{equation}
\delta E[A,\phase,N]=0.
\end{equation}
Stationarity alone is not enough: a maximum or saddle point does not define a stable particle. One therefore needs at least local non-negativity of the second variation within an allowed class of deformations,
\begin{equation}
\delta^2E[\eta]\geq 0,
\end{equation}
where $\eta$ denotes a small deformation of amplitude, phase, contour, or local profile.

The topological class acts as a restriction on allowed deformations. If the winding number $n$ is conserved, the node cannot continuously become vacuum without a discontinuity, profile suppression, or departure from the chosen state class:
\begin{equation}
n\neq 0 \quad \Rightarrow \quad \Gamma\not\sim\Gamma_{\mathrm{vac}}
\end{equation}
under fixed boundary conditions. This does not prove the existence of a particle, but it specifies the form of a future proof: one must find an action in which such classes exist and have finite energy.

A working test of the model should therefore proceed in this order:
\begin{enumerate}
\item choose the field space and boundary conditions;
\item construct an energy functional or action;
\item find stationary closed solutions;
\item check finite energy and absence of a point core;
\item check the spectrum of small perturbations around the solution;
\item compare stable classes with quantum numbers.
\end{enumerate}

Until these steps are performed, Field 01 remains a research program. The value of the program is that the hypothesis can be criticized in a concrete way: either the listed conditions can be realized, or they cannot.

\section{A Dynamical Core Instead of a Point Core}

The usual problem of a point-like particle is that if all structure is concentrated at a mathematical point, singular behavior may appear in the description. Field 01 replaces the point core by a dynamical core.

The core of the particle is not a point at which quantities must diverge. It is the closed self-consistency of the phase circulation. The particle is localized not because it contains a hard center, but because its phase rhythm returns to itself.

This allows the model to avoid an internal point singularity at the level of interpretation. It does not by itself solve all ultraviolet problems of quantum field theory. Rather, it suggests a different geometric intuition: elementary localization should be described by stable closed field modes rather than by structureless points.

\section{Memory as an Interphase Structure}

Field 01 uses the term memory not as psychological memory and not as a hidden classical record inside a particle. Memory means the retained pattern of phase relations that makes a configuration identifiable and dynamically continuous.

Let
\begin{equation}
\memory = \memory[\phase,\normal,\contour]
\end{equation}
denote the information retained by a closed configuration. This notation is schematic. It means that memory is carried not by a material shape alone, but by phase relations, circulation, and the stability of the local scalar / VEV-like profile.

This point is important for the later discussion of horizons. If memory is phase-structural rather than object-like, then a limiting surface can record the phase pattern without transporting a solid object through the horizon.

\section{Forces as Regimes of a Closed Wave}

The book-level formulation of Field 01 describes interactions in terms of how closed waves coordinate, hold, exchange, or open their phase structures. In this paper this idea is kept as a qualitative guide rather than as a completed derivation of the known interactions.

A cautious mapping is the following:
\begin{itemize}
\item electromagnetic interaction may be associated with phase asymmetry and open phase exchange;
\item strong interaction may be associated with tight coordination of local scalar profiles / VEV-profile layers and local phase locking;
\item weak interaction may be associated with instability, reconfiguration, or decay of a closed phase regime;
\item gravity may be associated with collective deformation of the field and with the energetic cost of stable configurations.
\end{itemize}

This mapping is not a substitute for gauge theory. The lines concerning strong and weak interaction require special caution: in established physics they are described by specific gauge symmetries and fields. Field 01 must either derive such structures as effective descriptions or clearly state the boundary of its applicability.

It is a research program: the next step would be to identify whether the phase variables of Field 01 can reproduce gauge symmetries, charges, and known conservation laws.

\section{Relation to Established Physics}

\subsection*{Quantum Field Theory}

Quantum field theory already treats particles as excitations of fields. Field 01 is compatible with this broad idea, but proposes a specific interpretation: a massive particle is not merely any excitation, but a closed and self-consistent phase regime.

The model should therefore be compared with solitons, topological defects, bound states, and other stable field configurations. The analogy is useful, but not yet a derivation.

\subsection*{Vacuum Correlations}

The vacuum in quantum field theory is not empty. It contains correlations and fluctuations. Field 01 uses this as a conceptual bridge: what appears as a particle may be treated as a stable pattern inside a deeper field structure.

No claim is made here that Field 01 has already derived the quantum vacuum. The claim is more modest: the language of phase structure and closure is not alien to modern field-based physics.

\subsection*{Spin}

Spin is a quantum number with precise transformation properties under rotations and Lorentz symmetry. Field 01 does not yet derive spin. It suggests that spin may eventually be related to the topology of phase circulation and to the way a closed wave returns to itself.

A future formalization would have to show how half-integer and integer spin representations arise. Until then, the relation between spin and closed circulation remains a hypothesis.

\subsection*{Higgs Mechanism}

The Standard Model explains elementary masses through coupling to the Higgs field. Field 01 does not replace this mechanism at the present stage. A possible bridge is to interpret mass as the energy of maintaining a stable closed configuration, while the Higgs mechanism supplies the established field-theoretic description of mass generation.

A serious comparison would require deriving effective mass terms from a Field 01 action and checking whether they can be related to known Higgs couplings.

\subsection*{Stress--Energy Tensor}

In general relativity the stress--energy tensor describes the distribution of energy and momentum that sources spacetime curvature. Field 01 may eventually require its own effective stress--energy tensor for phase configurations.

At the current stage, however, expressions such as phase tension or scalar-profile retention energy are only working terms. They should not be confused with an already defined tensorial theory.

\section{Limit: Vanishing Normal and Horizon}

The same logic that distinguishes an open photon mode from a closed massive particle also suggests a limiting regime near a horizon. If the local scalar / VEV-like profile is suppressed,
\begin{equation}
\normal \rightarrow 0,
\end{equation}
then the localized bulk-like structure can no longer be described as an ordinary closed particle in volume.

In Field 01 the horizon is interpreted as a recording surface. The closed configuration does not have to be imagined as mechanically torn apart. Rather, its phase memory can be transferred into a surface code:
\begin{equation}
\memory_{\mathrm{bulk}} \longrightarrow \memory_{\mathrm{boundary}}.
\end{equation}

This is intentionally close to the general spirit of the holographic principle and to the Bekenstein--Hawking relation between entropy and horizon area. It is not presented as a derivation of black-hole thermodynamics. It is a model-specific interpretation of why information may be associated with a boundary rather than with an internal object-like volume.

\section{Program for Further Formalization}

To turn Field 01 into a mathematical theory, several tasks are necessary.

\subsection*{State Space}

One must define the state space of the field:
\begin{equation}
\Psi_{01} \in \mathcal{H}_{01}
\end{equation}
or an equivalent classical field configuration space. It must be clear what the fundamental variables are and how physical states are identified.

\subsection*{Action or Energy Functional}

A dynamical principle is required. A possible schematic action is
\begin{equation}
S_{01} = \int \mathcal{L}(\phase, \partial_\mu \phase, \normal, \partial_\mu \normal)\, d^4x.
\end{equation}

The Lagrangian must be chosen or derived so that stable closed configurations exist.

\subsection*{Topological Charge}

Closed waves should be classified by a topological quantity, for example
\begin{equation}
Q = \frac{1}{2\pi}\oint_{\contour} d\phase.
\end{equation}

A stable particle-like solution would correspond to a non-trivial value of $Q$.

\subsection*{Stability Condition}

A particle should correspond to a stable or metastable extremum of the energy:
\begin{equation}
\delta E_{\mathrm{closed}} = 0, \qquad \delta^2 E_{\mathrm{closed}} > 0
\end{equation}
within an appropriate class of variations.

\subsection*{Bulk-to-Boundary Transition}

The transition toward a horizon should be described by a limiting map
\begin{equation}
\Pi_{\partial}: \memory_{\mathrm{bulk}} \rightarrow \memory_{\mathrm{boundary}},
\end{equation}
with the local scalar / VEV-like profile suppressed:
\begin{equation}
\normal \rightarrow 0.
\end{equation}

\subsection*{Comparison with Known Quantum Numbers}

Finally, the model must explain or reproduce known quantum numbers: charge, spin, mass, flavor, and gauge representation. Without this, Field 01 remains an interpretive framework rather than a predictive physical theory.

\section{What Is Claimed and What Remains Hypothetical}

The paper makes the following limited claims:
\begin{itemize}
\item it is possible to formulate a coherent phase-topological vocabulary for Field 01;
\item an elementary particle can be interpreted as a closed phase configuration;
\item the photon can be interpreted as open phase transport without a retained scalar / VEV-like profile;
\item mass can be associated with the energetic cost of maintaining a closed configuration;
\item a horizon can be interpreted as a limiting surface where the local scalar / VEV-like profile is suppressed and phase memory becomes boundary record.
\end{itemize}

The following points remain open:
\begin{itemize}
\item derivation from a precise action;
\item reproduction of the Standard Model particle spectrum;
\item derivation of spin and gauge charges;
\item quantitative relation to the Higgs mechanism;
\item construction of an effective stress--energy tensor;
\item testable predictions that differ from existing theories.
\end{itemize}

Thus Field 01 should currently be read as a structured research program, not as a completed theory of nature.

\section{Questions for Public Discussion}

For open technical discussion and later preprint work, the next mathematical iteration should answer questions such as:
\begin{enumerate}
\item Can one construct a simple energy functional that admits finite closed phase nodes?
\item Is one phase field enough, or is a multicomponent field, spinor structure, or gauge field required from the beginning?
\item Can the local scalar / VEV-like profile $N$ be defined without adding an unnecessary independent entity?
\item Do stable solutions with non-trivial winding and finite energy exist in three spatial dimensions?
\item Can twisted circulation be related to fermionic $4\pi$ periodicity without a mechanical rotation picture?
\item Can mass as the energy of a stable solution be compared with known mass terms and the Higgs mechanism?
\item What observable or theoretical distinction could separate Field 01 from a pure reformulation of known quantum field theory?
\end{enumerate}

These questions are deliberately formulated so that a negative answer is possible. If mathematical analysis shows that such closed phase nodes are impossible or unstable, that is important information for the model. If a minimal stable example is found, the next task is to compare its spectrum with established physics.

\section{Conclusion}

The central idea of this paper is simple: an elementary particle may be described as a closed wave of Field 01 rather than as a point-like object. The closure of phase circulation produces a localized node, the retained scalar / VEV-like profile gives the node local depth in the Field 01 interpretation, and the energy of maintaining this structure can be interpreted as mass.

A photon is then the limiting open case: it transports phase and energy, but does not retain a scalar / VEV-like profile and therefore has no rest mass. A horizon is the opposite limiting case: the local scalar / VEV-like profile is suppressed and the phase memory of a bulk configuration is transferred into a boundary record.

This language does not replace established physics. It provides a possible bridge between intuitive descriptions of particles, field-theoretic excitations, topological stability, and horizon information. The next task is to formulate the dynamics strictly enough that the model can be compared with known equations and, eventually, with observations.

\end{document}